\documentclass[11pt]{article}

% -------------------------------------------------------------------------------- %

% *** LANGUAGE ***
%\usepackage[italian]{babel}

% *** PAGE LAYOUT ***
\usepackage{geometry}
\geometry{a4paper,tmargin=2cm,bmargin=2cm,lmargin=2.5cm,rmargin=2.5cm}

% *** DEFAULT FONT SELECTION ***
\usepackage[T1]{fontenc}
\usepackage{lmodern}
\renewcommand*\familydefault{\sfdefault}

% *** GRAPHICS ***
\usepackage{graphicx}
\usepackage{xcolor}
\usepackage{tikz}
\usetikzlibrary{plotmarks}
\usepackage{pgfplots}
\pgfplotsset{compat=newest}
\pgfplotsset{plot coordinates/math parser=false}
\pgfplotsset{ grid style = {
    loosely dotted,
    line cap = round,
    black,
    line width = 0.5pt
  }
}
\usepackage[caption=false, font=footnotesize]{subfig}
\usepackage{float}

% *** AMS MATH ***
\usepackage{amsmath, amsfonts, amssymb}
\usepackage{empheq}
\usepackage{mathabx}

% *** TABLES ***
\usepackage{booktabs,topcapt}
\usepackage{multirow}

% *** TEXT FORMATTING ***
\usepackage[bottom]{footmisc}

% *** LISTS ***
\usepackage[shortlabels]{enumitem}

% *** LINE SPACING ***
\usepackage{setspace}

% *** PAGE NUMBERING ***
\usepackage{lastpage}

% *** SI UNITS ***
\usepackage{siunitx}

% *** BIBLIO ***
\usepackage{cite}

% *** MATLAB CODE ***
\usepackage{bigfoot}
\usepackage[numbered,framed]{matlab-prettifier}

\lstset{
backgroundcolor=\color[RGB]{247,247,247},
style=Matlab-editor,
basicstyle=\mlttfamily\scriptsize,
morekeywords={matlab2tikz},
keywordstyle=\color{blue},
morekeywords=[2]{1},
morekeywords={function},
keywordstyle=[2]{\color{black}},
identifierstyle=\color{black},
showstringspaces=false,
numbers=left,
numberstyle={\scriptsize \color[RGB]{147,147,147}},
numbersep=7pt,
aboveskip=16pt,
belowskip=16pt
}

\usepackage{verbatim}
\makeatletter
\def\verbatim@font{\ttfamily\scriptsize}
\makeatother

% *** OTHER ***
\usepackage{printlen}
\usepackage{tcolorbox}
\usepackage{import}

% *** EXTRA COMMANDS ***
\DeclareRobustCommand{\vect}[1]{ \boldsymbol{#1} }
\interfootnotelinepenalty=10000

% Notation shortcuts
\newcommand{\thh}{\vartheta_h}
\newcommand{\thb}{\vartheta_b}
\newcommand{\thd}{\vartheta_d}
\newcommand{\Jeq}{J_\mathrm{eq}}
\newcommand{\Beq}{B_\mathrm{eq}}
\newcommand{\wgc}{\omega_{gc}}
\newcommand{\wn}{\omega_n}

% -------------------------------------------------------------------------------- %

\title{
    {\Large Laboratory Report 3: Position Control of a DC Servomotor with Resonant Load}
}
\author{
    Name Lastname;\quad email: \texttt{\{name.lastname\}@unipd.it}\\[2pt]
    Group:~\texttt{[GROUP]},\quad Shift:~\texttt{[SHIFT]},\quad
    Members:~\texttt{[Lucio Farinati, Andrea Fasolo, Enrico Bassani, Francesco Bau]}
}

% -------------------------------------------------------------------------------- %

\usepackage{fancyhdr}
\lhead{} \chead{} \rhead{}
\lfoot{}
\cfoot{\textit{\scriptsize University of Padova}}
\rfoot{\scriptsize \thepage/\pageref{LastPage}}
\renewcommand{\headrulewidth}{0pt}
\renewcommand{\footrulewidth}{0pt}

% -------------------------------------------------------------------------------- %

\hyphenation{op-tical net-works semi-conduc-tor}

\begin{document}

\maketitle
\pagestyle{fancy}
\setcounter{page}{1}
\onehalfspacing

\sisetup{math-rm=\mathrm, text-rm=\rmfamily}

% ============================================================ %
\section{Introduction}
% ============================================================ %

\subsection{Activity Goal}

This report documents the work carried out in Laboratory~3 of the Control Engineering course.
The central objective is to synthesise angular-position controllers for a DC gearmotor whose
output shaft drives a compliant mechanical load — specifically, a rigid beam attached through
a torsional spring (elastic joint) rather than through a rigid coupling.
Such a configuration introduces a resonant mode into the plant dynamics, making the
control problem considerably more challenging than the single-inertia case addressed in
previous sessions.

Three controller families are investigated. A frequency-domain PID approach is applied first
because of its industrial prevalence and straightforward tuning via loop-shaping. Two
state-space alternatives are then explored: a pole-placement technique, which allows direct
assignment of the closed-loop eigenvalues, and an LQR-based optimisation strategy, which
trades off regulation performance against control effort through a quadratic cost criterion.

\subsection{System and Model}

\subsubsection*{Physical description}

% [TODO] Insert a labelled photograph or schematic of the experimental bench.
% At minimum, identify: (1) rigid beam, (2) elastic joint, (3) hub with potentiometer 2.

\begin{figure}[H]
    \centering
    % \includegraphics[width=0.62\textwidth]{figures/setup_lab3.jpg}
    \caption{Experimental bench: Quanser SRV-02 gearmotor fitted with the resonant-load
             attachment. Component labels --- (1) rigid beam, (2) pre-tensioned leaf-spring
             joint, (3) hub carrying potentiometer~2.}
    \label{fig:setup}
\end{figure}

The mechanical plant consists of a brushed DC motor (Faulhaber~2338S006S) connected to a
$N{=}14$ planetary gearbox, whose output shaft carries an aluminium hub of rotational
inertia~$J_h$.  A torsional spring (stiffness~$k$) couples the hub to an aluminium beam of
inertia~$J_b$, so that relative angular deflection $\thd = \thb - \thh$ can build up when
the hub accelerates.  Hub angular position~$\thh$ is sensed by an incremental quadrature
optical encoder ($500\,\text{ppr}$, ${\times}4$ decoding, $0.18\,\si{\degree/pulse}$).
Beam deflection~$\thd$ is converted to a voltage by a potentiometer mounted coaxially with
the joint pivot (potentiometer~2), subsequently digitised by a 16-bit ADC with
${\pm}10\,\si{V}$ range; the angular-to-voltage sensitivity is
$2.5\,\si{V}/170\,\si{\degree}$.

\subsubsection*{Mathematical model}

Neglecting the fast electrical and driver transients
($L_a/R_\mathrm{eq} \ll 1$, $T_\mathrm{drv} \ll 1$),
the translational-analogy equations for the two rotating bodies reduce to:
\begin{align}
    N^2 \Jeq\,\ddot{\thh} + N^2 \Beq\,\dot{\thh} + k\,\thh - k\,\thb
        &= \frac{N k_t k_\mathrm{drv}}{R_\mathrm{eq}}\,u - \tau_{d,h},
        \label{eq:hub}\\
    J_b\,\ddot{\thb} + B_b\,\dot{\thb} + k\,\thb - k\,\thh &= 0,
    \label{eq:beam}
\end{align}
where $\Jeq = J_m + J_h/N^2$ and $\Beq = B_m + B_h/N^2$ collect the inertia and viscous
damping referred to the load side, $u$~is the voltage command to the driver, and
$\tau_{d,h}$ lumps the Coulomb friction acting at the hub.  The restoring torque transmitted
through the elastic joint is $\tau_e = k\,(\thh - \thb)$.

Selecting $\vect{x}' = [\thh,\;\thd,\;\dot{\thh},\;\dot{\thd}]^T$ as state vector yields
the linear time-invariant model
\begin{equation}
    \dot{\vect{x}}' = A'\,\vect{x}' + B'\,u + B'_d\,\tau_{d,h}.
    \label{eq:ss}
\end{equation}
The explicit expressions for $A'$, $B'$, $B'_d$ are given in (52)--(54) of the theoretical
background document.  All symbols are defined in Table~\ref{tab:symbols}.

\begin{table}[H]
\centering
\caption{Notation for the resonant-load model.}
\label{tab:symbols}
\begin{tabular}{lll}
\toprule
Symbol & Meaning & Units \\
\midrule
$\thh$, $\thb$, $\thd$ & Hub, beam, and relative (deflection) angular positions & rad \\
$J_h$, $J_b$           & Hub and beam rotational inertia                         & \si{kg.m^2} \\
$B_h$, $B_b$           & Viscous damping at hub and beam                         & \si{N.m.s/rad} \\
$k$                    & Torsional stiffness of the elastic joint                 & \si{N.m/rad} \\
$\tau_{d,h}$           & Coulomb-friction disturbance torque at hub side          & \si{N.m} \\
$N$                    & Gearbox reduction ratio                                  & -- \\
$k_t$, $k_\mathrm{drv}$ & Motor torque constant, driver DC gain                  & \si{N.m/A}, -- \\
\bottomrule
\end{tabular}
\end{table}

\begin{table}[H]
\centering
\caption{Numerical values of the resonant-load parameters used in simulation.
         Starred entries (${}^*$) are initial guesses; Section~\ref{sec:estimation} reports
         the experimentally refined values.}
\label{tab:params}
\begin{tabular}{llll}
\toprule
Parameter & Symbol & Value & Source \\
\midrule
Hub inertia             & $J_h$         & $6.84\times10^{-4}$~\si{kg.m^2} & nominal (Al hub) \\
Hub viscous damping     & $B_h$         & from Lab~0                       & identified \\
Coulomb friction        & $\tau_{sf}$   & from Lab~0                       & identified \\
Beam inertia            & $J_b$         & $1.40\times10^{-3}$~\si{kg.m^2} & nominal \\
Beam viscous damping$^*$& $B_b$         & $3.4\times10^{-3}$~\si{N.m.s/rad} & initial guess \\
Joint stiffness$^*$     & $k$           & $0.83$~\si{N.m/rad}              & initial guess \\
\bottomrule
\end{tabular}
\end{table}

% ============================================================ %
\section{Tasks, Methodologies and Results}
% ============================================================ %

% ============================================================ %
\subsection{Simulink Model of the Two-Mass Drive (Task~2.1)}
% ============================================================ %

A discrete-time Simulink model was assembled to reproduce the closed-loop behaviour of the
complete electromechanical chain.  The model is structured into four communicating
subsystems: (i)~the \emph{drive}, which chains the DAC zero-order hold, the voltage
amplifier, the motor electrical dynamics, and the gearbox; (ii)~the \emph{mechanical
plant}, implementing equations~\eqref{eq:hub}--\eqref{eq:beam} in state-space form;
(iii)~the \emph{encoder}, performing quadrature counting and scaling; and (iv)~the
\emph{beam-deflection channel}, which applies the potentiometer gain and superimposes
zero-mean white Gaussian noise ($\sigma^2 = 3.5\times10^{-7}\,\si{V^2}$) followed by ADC
saturation and quantisation.  System-level screenshots appear in Appendix~\ref{app:simulink}.

All numerical constants were initialised through a single MATLAB script
(\texttt{load\_params\_resonant\_case.m}); the parameter set is reported in
Appendix~\ref{app:code}.  During the simulation phase (Tasks~2.2--2.4), the two unknowns
$B_b$ and $k$ were set to their nominal guesses (see Table~\ref{tab:params}); once
measured experimentally (Section~\ref{sec:estimation}), the same model was re-run with the
updated values.

% ============================================================ %
\subsection{PID Position Controller (Tasks~2.2 and 3.3)}
% ============================================================ %

\subsubsection*{Frequency-domain synthesis}

A PID with real derivative filter was adopted to regulate the collocated output~$\thh$:
\begin{equation}
    C(s) = K_P\!\left(1 + \frac{1}{T_I s} + \frac{T_D\,s}{T_L\,s + 1}\right),
    \label{eq:pid}
\end{equation}
with the derivative filter bandwidth chosen as $1/T_L \approx 10\,\wgc$ to attenuate
high-frequency noise while preserving adequate phase contribution near crossover.
An integrator anti-windup back-calculation scheme was included; the reset time constant
was set to $T_W = t_{s,5\%}/5$ (anti-windup gain $K_W = 1/T_W$).

The desired closed-loop transient — settling time $t_{s,5\%} \leq 0.85\,\si{s}$ and
overshoot $M_p \leq 30\,\%$ — fixes a target gain crossover frequency $\wgc$ and
a minimum phase margin.  Under the constraint $T_I/T_D = 4$, the three free parameters
$K_P$, $T_I$, $T_D$ were obtained by imposing the desired magnitude and phase conditions
at $\wgc$ on the loop transfer function $L(s) = C(s)\,P_{u\to\thh}(s)$.
The resulting values are collected in Table~\ref{tab:pid_params_sim}.

\begin{table}[H]
\centering
\caption{PID parameters obtained during the simulation phase (tentative $k$, $B_b$).}
\label{tab:pid_params_sim}
\begin{tabular}{cccccc}
\toprule
$K_P$ & $T_I$~[s] & $T_D$~[s] & $T_L$~[s] & $T_W$~[s] & $\wgc$~[rad/s] \\
\midrule
\texttt{[---]} & \texttt{[---]} & \texttt{[---]} & \texttt{[---]} & \texttt{[---]} & \texttt{[---]} \\
\bottomrule
\end{tabular}
\end{table}

\subsubsection*{Simulation results}

% [TODO] Replace placeholder figures with your actual plots.
% Each figure should show: time axis [s], angle axis [deg], reference line, ts,5% band.
% Use solid/dashed to distinguish with/without AW; add a legend.

\begin{figure}[H]
    \centering
    % \includegraphics[width=0.90\textwidth]{figures/pid_aw_50deg_sim.pdf}
    \caption{Simulated hub-position response to a $50\,\si{\degree}$ step command.
             Solid line: PID with back-calculation anti-windup active.
             Dashed line: integrator unclamped (no anti-windup).
             Tentative elastic-joint parameters.}
    \label{fig:pid_aw_50_sim}
\end{figure}

\begin{figure}[H]
    \centering
    % \includegraphics[width=0.90\textwidth]{figures/pid_aw_120deg_sim.pdf}
    \caption{Same controller, $120\,\si{\degree}$ step command. The divergence between the
             two curves grows with reference amplitude because a larger step drives the
             integrator further into saturation, making windup more severe.}
    \label{fig:pid_aw_120_sim}
\end{figure}

\begin{table}[H]
\centering
\caption{Transient performance indices from simulation: effect of the anti-windup mechanism.}
\label{tab:pid_perf_sim}
\begin{tabular}{llcc}
\toprule
Reference & Anti-windup & $t_{s,5\%}$~[s] & $M_p$~[\%] \\
\midrule
$50\,\si{\degree}$  & enabled  & \texttt{[---]} & \texttt{[---]} \\
$50\,\si{\degree}$  & disabled & \texttt{[---]} & \texttt{[---]} \\
$120\,\si{\degree}$ & enabled  & \texttt{[---]} & \texttt{[---]} \\
$120\,\si{\degree}$ & disabled & \texttt{[---]} & \texttt{[---]} \\
\bottomrule
\end{tabular}
\end{table}

% [TODO] Comment: quantify the improvement brought by anti-windup for the large step.
% Explain the mechanism: during saturation the integrator accumulates excess error;
% without clamping, this produces a large overshoot upon exit from saturation.

\subsubsection*{Redesign with measured joint parameters and experimental validation}

Once $\hat{k}$ and $\hat{B}_b$ became available from Section~\ref{sec:estimation}, the
loop-shaping procedure was repeated on the updated plant model.  The revised parameters
are listed in Table~\ref{tab:pid_params_exp}; Figure~\ref{fig:pid_exp_vs_sim} overlays
the corresponding simulation trace with the response acquired on the physical bench.

\begin{table}[H]
\centering
\caption{Revised PID parameters after substituting experimentally identified $\hat{k}$
         and $\hat{B}_b$.}
\label{tab:pid_params_exp}
\begin{tabular}{cccccc}
\toprule
$K_P$ & $T_I$~[s] & $T_D$~[s] & $T_L$~[s] & $T_W$~[s] & $\wgc$~[rad/s] \\
\midrule
\texttt{[---]} & \texttt{[---]} & \texttt{[---]} & \texttt{[---]} & \texttt{[---]} & \texttt{[---]} \\
\bottomrule
\end{tabular}
\end{table}

\begin{figure}[H]
    \centering
    % \includegraphics[width=0.90\textwidth]{figures/pid_exp_vs_sim.pdf}
    \caption{Hub-position step response ($50\,\si{\degree}$): numerical simulation with
             identified parameters (dashed) against physical-plant measurement (solid),
             revised PID with anti-windup.}
    \label{fig:pid_exp_vs_sim}
\end{figure}

\begin{table}[H]
\centering
\caption{PID closed-loop performance: numerical model vs.\ physical experiment.}
\label{tab:pid_perf_exp}
\begin{tabular}{lcc}
\toprule
& $t_{s,5\%}$~[s] & $M_p$~[\%] \\
\midrule
Simulation   & \texttt{[---]} & \texttt{[---]} \\
Experimental & \texttt{[---]} & \texttt{[---]} \\
\bottomrule
\end{tabular}
\end{table}

% [TODO] Discuss residual discrepancies (nonlinear friction, model order reduction, etc.).

% ============================================================ %
\subsection{Elastic Joint Parameter Identification (Task~3.2)}
\label{sec:estimation}
% ============================================================ %

With the hub shaft locked by a mechanical stop, the beam executes a damped oscillation
governed solely by its own inertia $J_b$, its viscous damping $B_b$, and the joint
restoring torque.  The deflection angle therefore satisfies
\begin{equation}
    J_b\,\ddot{\thd} + B_b\,\dot{\thd} + k\,\thd = 0,
    \label{eq:free_osc}
\end{equation}
a second-order underdamped oscillator with undamped natural frequency
$\wn = \sqrt{k/J_b}$ and damping ratio $\delta = B_b/(2\sqrt{J_b\,k})$.

A free-oscillation test was performed by manually deflecting the beam to
$\thd(0) \approx \texttt{[---]}\,\si{\degree}$ (kept below $20\,\si{\degree}$ to remain
within the linear spring regime) and releasing it.  The decaying signal was recorded at
sampling period $T_s = 1\,\si{ms}$.

The two unknown parameters were extracted by the logarithmic-decrement method: the
$M$~successive peak magnitudes $\{|\thd(t_k)|\}$ satisfy
$\log|\thd(t_k)| = -k\,\xi + b$, where $\xi = \delta\pi/\sqrt{1-\delta^2}$ is the
logarithmic decrement.  A least-squares line fit on the
$(k,\,\log|\thd(t_k)|)$ plane yields $\hat{\xi}$ and, through the inverse relation
$\hat{\delta} = \hat{\xi}/\sqrt{\pi^2 + \hat{\xi}^2}$, the damping ratio.
The oscillation frequency~$\hat{\omega}$ was averaged across all successive
peak-to-peak intervals; the natural frequency follows as
$\hat{\wn} = \hat{\omega}/\sqrt{1-\hat{\delta}^2}$.

The final estimates and the derived physical parameters are:
\begin{equation}
    \hat{\wn} = \texttt{[---]}\;\si{rad/s},\quad
    \hat{\delta} = \texttt{[---]},\quad
    \hat{k} = J_b\,\hat{\wn}^2 = \texttt{[---]}\;\si{N.m/rad},\quad
    \hat{B}_b = 2\,J_b\,\hat{\delta}\,\hat{\wn} = \texttt{[---]}\;\si{N.m.s/rad}.
    \label{eq:estimates}
\end{equation}

\begin{figure}[H]
    \centering
    % \includegraphics[width=0.85\textwidth]{figures/free_response_fit.pdf}
    \caption{Recorded free-oscillation signal $\thd(t)$ with the hub locked (solid line),
             overlaid with the fitted exponential amplitude envelope
             $\pm\hat{A}\,e^{-\hat{\delta}\hat{\wn}\,t}$ (dashed lines).}
    \label{fig:free_resp}
\end{figure}

% ============================================================ %
\subsection{State-Space Control via Eigenvalue Assignment (Tasks~2.3 and 3.4)}
% ============================================================ %

\subsubsection*{Pole-placement synthesis}

For state-space design the plant is represented by~\eqref{eq:ss} with the four-component
state $\vect{x}'$.  Hub and beam velocities are not directly measured; they are
reconstructed by filtering the position signals through a second-order high-pass network:
\begin{equation}
    H_1(s) = \frac{\wc^2\,s}{s^2 + \sqrt{2}\,\wc\,s + \wc^2},
    \qquad \wc = 2\pi\cdot 50\;\si{rad/s},
    \label{eq:hpf}
\end{equation}
which approximates differentiation below $\wc$ while limiting noise amplification above it.

The desired closed-loop behaviour ($t_{s,5\%} \le 0.85\,\si{s}$, $M_p \le 30\,\%$)
maps to a dominant second-order pair with
\begin{equation}
    \wn = \frac{3}{\delta\,t_{s,5\%}} = \texttt{[---]}\;\si{rad/s},\qquad
    \delta = \frac{\ln(1/M_p)}{\sqrt{\pi^2 + \ln^2(1/M_p)}} = \texttt{[---]}.
    \label{eq:wn_delta}
\end{equation}
The auxiliary angle $\phi = \arctan\!\bigl(\sqrt{1-\delta^2}/\delta\bigr)$ defines the
target sector.  Since the model order is four, two additional eigenvalues are placed at
$\omega_n\,e^{j(-\pi\pm\phi/2)}$, which lie further to the left and therefore decay faster
than the dominant pair:
\begin{equation}
    \lambda_{c,\{1,2\}} = \wn\,e^{j(-\pi\pm\phi)},\qquad
    \lambda_{c,\{3,4\}} = \wn\,e^{j(-\pi\pm\phi/2)}.
    \label{eq:eigenvalues}
\end{equation}

The state feedback vector $\vect{K}\in\mathbb{R}^{1\times4}$ was computed via the MATLAB
\texttt{place} function applied to the pair $(A',\,B')$:
\begin{equation}
    \vect{K} = \bigl[\;\texttt{[---]}\;\;\texttt{[---]}\;\;\texttt{[---]}\;\;\texttt{[---]}\;\bigr].
\end{equation}

\subsubsection*{Numerical and experimental results}

\begin{figure}[H]
    \centering
    % \includegraphics[width=0.90\textwidth]{figures/ss_eigen_50deg.pdf}
    \caption{Hub-position step response ($50\,\si{\degree}$) with the pole-placement
             regulator (nominal tracking): simulation (dashed) vs.\ physical measurement
             (solid).  Identified joint parameters.}
    \label{fig:ss_eigen_50}
\end{figure}

\begin{table}[H]
\centering
\caption{Pole-placement controller: transient performance, simulation vs.\ experiment.}
\label{tab:ss_eigen_perf}
\begin{tabular}{lcc}
\toprule
& $t_{s,5\%}$~[s] & $M_p$~[\%] \\
\midrule
Simulation   & \texttt{[---]} & \texttt{[---]} \\
Experimental & \texttt{[---]} & \texttt{[---]} \\
\bottomrule
\end{tabular}
\end{table}

% [TODO] Optional: report the integral-action extension (Tasks 2.3.3-4)
%        for robust rejection of constant load torques.

% ============================================================ %
\subsection{State-Space Control via LQR Optimisation (Tasks~2.4 and 3.5)}
% ============================================================ %

\subsubsection*{Design based on the Symmetric Root Locus (Task~2.4.1)}

The first LQR formulation penalises only the hub-position error and the control voltage:
\begin{equation}
    J = \int_0^{+\infty} \thh^2(t) + r\,u^2(t)\;\mathrm{d}t, \qquad r > 0.
    \label{eq:cost_srl}
\end{equation}
Rather than guessing $r$ by trial and error, the scalar weight was chosen graphically
by plotting the Symmetric Root Locus (SRL) of the loop $G(-s)G(s)$, where $G(s)$ is the
transfer function from $u$ to $\thh$ of the reduced model~\eqref{eq:ss}.  The SRL reveals
how the optimal closed-loop eigenvalues migrate as $1/r$ increases from zero; the admissible
region in the left-half plane — $\mathrm{Re}(\lambda) \le -3/t_{s,5\%}$ and phase angle
within the sector bounded by $\pm\phi$ — provides a geometric selection criterion.

\begin{figure}[H]
    \centering
    % \subfloat[Full SRL]{%
    %   \includegraphics[width=0.47\textwidth]{figures/srl_full.pdf}}
    % \hfill
    % \subfloat[Selected eigenvalue locations]{%
    %   \includegraphics[width=0.47\textwidth]{figures/srl_selected.pdf}}
    \caption{Symmetric Root Locus for the cost function~\eqref{eq:cost_srl}.
             Left panel: full locus with admissible region (shaded). Right panel:
             closed-loop eigenvalue locations (squares) corresponding to the chosen
             $r = \texttt{[---]}$.}
    \label{fig:srl}
\end{figure}

The selected weight $r = \texttt{[---]}$ leads to
\begin{equation}
    \vect{K}_\mathrm{SRL} = \bigl[\;\texttt{[---]}\;\;\texttt{[---]}\;\;\texttt{[---]}\;\;\texttt{[---]}\;\bigr].
\end{equation}

\subsubsection*{Design based on Bryson's normalisation rule (Task~2.4.3)}

For the second LQR variant the cost is extended to all state components:
\begin{equation}
    J = \int_0^{+\infty} \vect{x}^{\prime T} Q\,\vect{x}' + r\,u^2\;\mathrm{d}t.
    \label{eq:cost_bryson}
\end{equation}
The diagonal weight matrix $Q$ was initialised by Bryson's normalisation: each entry
$q_{ii} = 1/\bar{x}_i^2$, where $\bar{x}_i$ is the largest acceptable deviation of the
$i$-th state component.  For a $50\,\si{\degree}$ reference the tolerances are set as
$\bar{\thh} = 0.3 \times 50 \times (\pi/180)$, $\bar{\thd} = \pi/36$, and the
velocity weights were left at zero (fast velocity errors are inherently damped).
The input weight follows as $r = 1/\bar{u}^2$ with $\bar{u} = 10\,\si{V}$:
\begin{equation}
    Q = \mathrm{diag}\!\left(\frac{1}{\bar{\thh}^2},\;\frac{1}{\bar{\thd}^2},\;0,\;0\right),
    \qquad r = 10^{-2}.
    \label{eq:bryson_weights}
\end{equation}
The optimal feedback vector is
\begin{equation}
    \vect{K}_\mathrm{Bry} = \bigl[\;\texttt{[---]}\;\;\texttt{[---]}\;\;\texttt{[---]}\;\;\texttt{[---]}\;\bigr].
\end{equation}

\subsubsection*{Simulation and experimental comparison}

\begin{figure}[H]
    \centering
    % \includegraphics[width=0.90\textwidth]{figures/lqr_compare_50deg.pdf}
    \caption{Hub-position step response ($50\,\si{\degree}$): SRL-based LQR (blue) and
             Bryson-normalised LQR (red) under simulation (dashed) and physical test
             (solid).  Both runs use identified joint parameters.}
    \label{fig:lqr_compare}
\end{figure}

\begin{table}[H]
\centering
\caption{LQR step-response performance ($50\,\si{\degree}$): both designs, simulation
         vs.\ experiment.}
\label{tab:lqr_perf}
\begin{tabular}{llcc}
\toprule
Controller & & $t_{s,5\%}$~[s] & $M_p$~[\%] \\
\midrule
$\vect{K}_\mathrm{SRL}$  & Simulation   & \texttt{[---]} & \texttt{[---]} \\
                          & Experimental & \texttt{[---]} & \texttt{[---]} \\
\midrule
$\vect{K}_\mathrm{Bry}$  & Simulation   & \texttt{[---]} & \texttt{[---]} \\
                          & Experimental & \texttt{[---]} & \texttt{[---]} \\
\bottomrule
\end{tabular}
\end{table}

% [TODO] Compare the two LQR designs: discuss differences in settling time,
%        overshoot, and beam-deflection amplitude (thd).
%        Notably: does penalising thd in Bryson's Q reduce resonant oscillations?
% [TODO] If implemented, add FS-LQR results (Tasks 2.4.9-10) here or as a separate subsection.

% ============================================================ %
\section*{Summary and Conclusions}
% ============================================================ %

% [TODO] 3-6 sentences synthesising the key outcomes:
%   - Which controllers met the specs? By how much?
%   - Main gap between simulation and physical experiment, and likely causes.
%   - Comparative assessment of PID vs. pole placement vs. LQR for this plant.

% ============================================================ %
\appendix

\section{MATLAB Code}
\label{app:code}

\begin{lstlisting}
\end{lstlisting}

\begin{lstlisting}
\end{lstlisting}

\begin{lstlisting}
\end{lstlisting}

% -------------------------------------------------------------------------------- %
\appendix
\section{Appendix}
\subsection{Matlab code}
\subsection{Data sheet}
\subsection{Detailed calculations for PID design}

\label{LastPage}
\end{document}
